\subsection{Ample comparison quotients on Hirzebruch surfaces}
\label{subsec:hirzebruch-ample-comparison}

Let $e\geq0$ be an integer and let
\[
 X=\mathbb P_{\mathbb P^1}(\mathcal O\oplus\mathcal O(e)),
 \qquad \pi:X\longrightarrow\mathbb P^1,
\]
where the projective bundle parametrizes lines. Write $S$ for the section defined by the subbundle $\mathcal O(e)$ and $f$ for a ruling fibre. The normal bundle of $S$ is $\mathcal O(-e)$, so
\[
 S^2=-e,\qquad S\cdot f=1,\qquad f^2=0.
\]
The tautological dual is $\mathcal O_X(S)$. A line bundle with fibre degree zero is pulled back from $\mathbb P^1$: its direct image has rank one and the fibrewise evaluation is an isomorphism. Consequently $\operatorname{Pic}(X)=\mathbb ZS\oplus\mathbb Zf$.

\begin{lemma}\label{lem:hirzebruch-ruling-sections}
For integers $a\geq0$ and $b$,
\[
 \pi_*\mathcal O_X(aS+bf)=
 \bigoplus_{i=0}^a\mathcal O_{\mathbb P^1}(b-ie),
 \qquad R^1\pi_*\mathcal O_X(aS+bf)=0.
\]
A class $aS+bf$ is effective exactly when $a,b\geq0$. If $a\geq1$ and $b>ae$, its line bundle is very ample. In particular,
\[
 H=S+(3e+1)f,
 \qquad H^2=5e+2,
 \qquad h^0(X,\mathcal O_X(H))=5e+4.
\]
The canonical class is $K_X=-2S-(e+2)f$.
\end{lemma}
\begin{proof}
On a trivializing base chart, sections of the tautological dual power are homogeneous fibre polynomials of degree $a$. Their transition law gives
\[
 \pi_*\mathcal O_X(aS)=
 \operatorname{Sym}^a(\mathcal O\oplus\mathcal O(-e)).
\]
Twisting on the base gives the formula. The first fibre cohomology of $\mathcal O_{\mathbb P^1}(a)$ vanishes, giving the higher-direct-image assertion. If $a<0$, restriction to a general ruling fibre excludes a nonzero section. For $a\geq0$, the direct-image formula gives a nonzero section exactly when $b\geq0$.

Suppose $a\geq1$ and $b>ae$. All coefficient line bundles in the displayed direct image have degree at least one. Fibre polynomials of degree $a$ separate points and tangent directions within a fibre. To separate points over different base points, choose a fibre monomial nonzero at the first point and a coefficient section vanishing at the second base point but not the first. At any point, choose such a nonzero monomial and a coefficient section with a simple zero at its base point. The resulting global section vanishes there, has nonzero derivative in a prescribed base direction and has zero vertical derivative. The complete linear system therefore separates points and first jets, so it defines an embedding. The formulas for $H^2$ and $h^0(H)$ follow from the intersections and the two summands of its direct image. The relative Euler sequence gives $K_{X/\mathbb P^1}=-2S-ef$; adding the base canonical class proves the last identity.
\end{proof}

Put
\[
 L_1=\mathcal O_X(5S-f),\qquad
 L_2=\mathcal O_X(-S+(12e+5)f),\qquad
 E_0=L_1\oplus L_2.
\]
Their Chern classes satisfy
\[
 c_1(E_0)=4H,\qquad c_2(E_0)=65e+26=13H^2.
\]
An ample bundle with these Chern classes can be constructed as a quotient of a trivial bundle. The essential point is to exclude common curves uniformly over the dual projective parameter.

\begin{proposition}\label{prop:hirzebruch-ample-comparison}
There exist linearly independent sections $a_0,\ldots,a_3$ of $\mathcal O_X(H)$ with no common zero. For any such quadruple and a general quadruple $b_0,\ldots,b_3$ of sections of $\mathcal O_X(3H)$, the sequence
\[
 0\longrightarrow\mathcal O_X(-H)\oplus\mathcal O_X(-3H)
 \xrightarrow{(a,b)}\mathcal O_X^{\oplus4}
 \longrightarrow E_{\mathrm{amp}}\longrightarrow0
\]
has a locally free ample cokernel of rank two, with $c_1(E_{\mathrm{amp}})=4H$ and $c_2(E_{\mathrm{amp}})=13H^2$. Here general means outside a proper Zariski-closed subset of the space of quadruples.
\end{proposition}
\begin{proof}
Global generation makes a common zero at a fixed point impose four independent linear conditions on a quadruple $a$. The incidence with a moving point of the surface has codimension at least two in the space of quadruples; its projection is closed by projectivity. Independence is also a nonempty open condition because $h^0(H)\geq4$. The two open conditions intersect. Fix such an $a$ and write
\[
 V=H^0(X,\mathcal O_X(3H))^{\oplus4},\qquad N=\dim V.
\]
At a fixed $x$, rank drops exactly when $b(x)$ is proportional to the nonzero $a(x)$. Surjective evaluation makes this a codimension-three condition. Its incidence over $X$ has dimension at most $N-1$ and closed image in $V$. Outside that image the cokernel is locally free of rank two.

The dual kernel is a subbundle of $(\mathbb C^4)^*\otimes\mathcal O_X$. It defines a proper morphism
\[
 q_b:\mathbb P(E_{\mathrm{amp}}^*)\longrightarrow
 \mathbb P((\mathbb C^4)^*),
 \qquad \mathcal O_{\mathbb P(E_{\mathrm{amp}}^*)}(1)
 =q_b^*\mathcal O_{\mathbb P^3}(1).
\]
For $w=[w_0:\cdots:w_3]$, its fibre is the common zero scheme of
\[
 s_w=\sum_iw_i a_i\in H^0(X,\mathcal O_X(H)),
 \qquad t_w=\sum_iw_i b_i\in H^0(X,\mathcal O_X(3H)).
\]
Independence of $a$ gives $s_w\neq0$ for every $w$.

A nonzero effective subdivisor $D$ of $\operatorname{div}(s_w)$ has class $uS+vf$, where
\[
 u\in\{0,1\},\qquad 0\leq v\leq3e+1,
 \qquad (u,v)\neq(0,0).
\]
Write $|D|=\mathbb P H^0(X,\mathcal O_X(D))$ for the complete linear system. For each of these finitely many classes $D_0$, let $T_{D_0}$ parametrize triples $(w,[\sigma],[\tau])$ with
\[
 [\sigma]\in|D_0|,\qquad [\tau]\in|H-D_0|,
 \qquad [s_w]=[\sigma\tau].
\]
Multiplication of nonzero sections is nonzero, so this is a closed projective parameter scheme. Its fibre over $w$ consists of the finitely many effective subdivisors of the fixed divisor $\operatorname{div}(s_w)$, with their possible multiplicities. Thus $\dim T_{D_0}\leq3$.

Vanishing on any nonzero effective divisor $D$ imposes at least four independent conditions on $H^0(X,\mathcal O_X(3H))$. Indeed, choose an integral curve $C$ in its support. Very ampleness supplies sections $r_0,r_1$ of $\mathcal O_X(H)$ with nonconstant ratio on $C$. The four cubic monomials in their restrictions are independent: a relation would make the nonconstant ratio algebraic over $\mathbb C$, hence constant.

For fixed $(w,[\sigma],[\tau])$, the map $V\to H^0(X,\mathcal O_X(3H))$, $b\mapsto t_w$, is surjective. The condition that $t_w$ vanish on $D=\operatorname{div}(\sigma)$ therefore has codimension at least four. It is closed in families: the universal multiplication map
\[
 H^0(X,\mathcal O_X(3H-D_0))\otimes\mathcal O_{|D_0|}(-1)
 \longrightarrow H^0(X,\mathcal O_X(3H))\otimes\mathcal O_{|D_0|}
\]
is fibrewise injective, so its image is a subbundle. Vanishing of the image of $t_w$ in its quotient defines the condition. Its incidence over $T_{D_0}$ has dimension at most $N-1$ and closed image in $V$.

Choose $b$ outside these finitely many proper closed images and the rank-drop image. The sections $s_w,t_w$ then have no common curve for any $w$, so every fibre of $q_b$ is finite. The proper morphism is finite, and its pullback of $\mathcal O_{\mathbb P^3}(1)$ is ample. Finally,
\[
 c(E_{\mathrm{amp}})=\frac1{(1-H)(1-3H)}=1+4H+13H^2.
\]
\end{proof}

A vector bundle $E$ on $X$ is called \emph{prioritary for the ruling} when $\operatorname{Ext}^2_X(E,E(-f))=0$. All ruling fibres are linearly equivalent. Serre duality identifies this condition with
\[
 \operatorname{Hom}_X(E,E(K_X+f))=0.
\]

\begin{lemma}\label{lem:hirzebruch-prioritary-comparison}
Both $E_0$ and $E_{\mathrm{amp}}$ are prioritary for the ruling.
\end{lemma}
\begin{proof}
Write $D=K_X+f=-2S-(e+1)f$. The summands of $\operatorname{End}(E_0)(D)$ have classes
\[
 D,\quad D,\quad 4S-(13e+7)f,\quad -8S+(11e+5)f.
\]
Each has a negative coefficient in the effective-cone basis, so each has zero space of sections. This proves the required Hom vanishing for $E_0$.

Precomposition with $\mathcal O_X^4\twoheadrightarrow E_{\mathrm{amp}}$ injects $\operatorname{Hom}(E_{\mathrm{amp}},E_{\mathrm{amp}}(D))$ into $H^0(X,E_{\mathrm{amp}}(D))^4$. Twist the defining sequence by $D$. Its middle term has no sections. For $j=1,3$, Serre duality and the ruling formula give
\[
 H^1(X,\mathcal O_X(D-jH))
 =H^1(X,\mathcal O_X(jH-f))^*,
 \qquad
 \pi_*\mathcal O_X(jH-f)
 =\bigoplus_{i=0}^j\mathcal O_{\mathbb P^1}(j(3e+1)-1-ie).
\]
The higher direct image is zero, and every displayed degree is at least $j(2e+1)-1\geq0$. Hence all these first cohomology groups vanish. The twisted exact sequence gives $H^0(X,E_{\mathrm{amp}}(D))=0$, proving the second Hom vanishing.
\end{proof}

\begin{proposition}\label{prop:hirzebruch-positive-dual-tests}
For $F_s=L_s^*$ there are K\"ahler forms $\eta_s$ with $\deg_{\eta_s}F_s=1$. Consequently the split bundle $E_0^*=F_1\oplus F_2$ satisfies the hypotheses of \cref{thm:convex-norm-deformation-obstruction}.
\end{proposition}
\begin{proof}
The integral classes
\[
 A_1=6S+(6e+1)f,\qquad A_2=S+(13e+6)f
\]
are very ample by \cref{lem:hirzebruch-ruling-sections}. Pullbacks of normalized Fubini--Study forms under their embeddings give K\"ahler representatives $\eta_s$. Direct intersection gives
\[
 c_1(L_1)\cdot A_1=-1,\qquad
 c_1(L_2)\cdot A_2=-1.
\]
Dualizing changes each degree to one. The closed positive $(1,1)$-forms $\eta_s$ are the required individual Gauduchon tests in dimension two.
\end{proof}
